\chapter*{Anexos}
\addcontentsline{toc}{chapter}{Anexos}

\section*{Anexo A: Instrumento de Consenso Delphi}

El cuestionario siguiente fue aplicado en dos rondas Delphi a un panel de 12
expertos del sector bancario peruano (arquitectos de software, líderes
técnicos y oficiales de cumplimiento). La primera ronda consensuó los
criterios técnicos y regulatorios clásicos; la segunda ronda incorporó
específicamente los dos criterios emergentes asociados a la generación de
código asistida por Inteligencia Artificial. El objetivo fue consensuar los
criterios de evaluación y sus pesos relativos para la selección de
frameworks de interfaz de usuario en instituciones financieras reguladas.

\subsection*{Escala de valoración}

\begin{center}
\begin{tabular}{@{}cp{10cm}@{}}
\toprule
Valor & Significado \\
\midrule
1 & Sin relevancia para el entorno bancario \\
2 & Relevancia baja \\
3 & Relevancia moderada \\
4 & Relevancia alta \\
5 & Relevancia crítica --- factor determinante en la selección \\
\bottomrule
\end{tabular}
\end{center}

\subsection*{Sección I --- Criterios técnicos}

\begin{center}
\small
\begin{tabular}{@{}cp{2.8cm}p{8cm}c@{}}
\toprule
N.\textsuperscript{o} & Criterio & Descripción para el contexto bancario & Valor (1--5) \\
\midrule
1 & Seguridad intrínseca & Mecanismos nativos contra XSS, CSRF y vulnerabilidades OWASP Top 10 sin configuración adicional & \\
2 & Rendimiento de carga & Métricas Core Web Vitals (LCP $\leq$ 2.5\,s, INP $\leq$ 200\,ms, CLS $\leq$ 0.1) en canales digitales con tráfico pico & \\
3 & Mantenibilidad & Facilidad de refactorización, cobertura de pruebas automatizadas y ciclos de actualización & \\
4 & Calidad del código IA & Consistencia arquitectónica y auditabilidad de las sugerencias de asistentes de código (Copilot, Cursor) & \\
5 & Accesibilidad & Conformidad nativa con WCAG 2.2 nivel AA sin librerías de terceros & \\
6 & Compatibilidad multiplataforma & Consistencia cross-browser y omnicanal (web, móvil, cajeros) sin degradación de funcionalidad & \\
\bottomrule
\end{tabular}
\end{center}

\subsection*{Sección II --- Criterios normativos y de cumplimiento}

\begin{center}
\small
\begin{tabular}{@{}cp{2.8cm}p{8cm}c@{}}
\toprule
N.\textsuperscript{o} & Criterio & Descripción para el contexto bancario & Valor (1--5) \\
\midrule
7 & Cumplimiento PCI-DSS & Facilidad de implementar controles del estándar PCI-DSS v4.0 en la capa de presentación & \\
8 & Integración legacy & Adaptabilidad a arquitecturas core bancarias de 15+ años de antigüedad & \\
9 & SLA y disponibilidad & Capacidad para sostener disponibilidad $\geq$ 99.9\% con picos de tráfico de 15$\times$ la línea base & \\
10 & Gobernanza de IA & Trazabilidad, riesgo de propiedad intelectual y auditabilidad regulatoria del código generado por asistentes de IA & \\
\bottomrule
\end{tabular}
\end{center}

\subsection*{Criterios adicionales sugeridos por el experto}

\noindent\hrulefill

\vspace{0.3cm}
\noindent\hrulefill

\subsection*{Perfil del experto}

\begin{center}
\begin{tabular}{@{}p{5cm}p{7cm}@{}}
\toprule
Años de experiencia en banca: & \makebox[7cm]{\hrulefill} \\[0.4cm]
Rol actual: & \makebox[7cm]{\hrulefill} \\[0.4cm]
Frameworks de interfaz de usuario que conoce: & \makebox[7cm]{\hrulefill} \\[0.4cm]
¿Utiliza asistentes de código basados en IA? & \makebox[7cm]{\hrulefill} \\[0.4cm]
Institución (opcional): & \makebox[7cm]{\hrulefill} \\
\bottomrule
\end{tabular}
\end{center}

\section*{Anexo B: Pesos AHP y Evaluación Detallada del Caso de Estudio}

\subsection*{B.1 Vector de prioridades derivado del proceso AHP}

La Tabla \ref{tab:4.6} presenta los pesos normalizados obtenidos tras aplicar
el Proceso Analítico Jerárquico (AHP) con la participación del panel de
expertos, incluida la segunda ronda Delphi que incorporó los criterios de IA
generativa. La razón de consistencia calculada fue $CR = 0.06 \leq 0.10$,
confirmando la validez de los juicios de comparación.

\begin{table}[htbp]
\centering
\caption{Vector de prioridades AHP --- Caso de estudio banco peruano mediano}
\label{tab:4.6}
\small
\begin{tabular}{@{}clcc@{}}
\toprule
N.\textsuperscript{o} & Criterio ($C_i$) & Peso $w_i$ & Porcentaje \\
\midrule
1  & Seguridad intrínseca          & 0.16 & 16.0\% \\
2  & Rendimiento de carga          & 0.13 & 13.0\% \\
3  & Mantenibilidad                & 0.13 & 13.0\% \\
4  & Calidad del código IA         & 0.08 & 8.0\%  \\
5  & Accesibilidad WCAG            & 0.06 & 6.0\%  \\
6  & Compatibilidad cross-browser  & 0.04 & 4.0\%  \\
7  & Cumplimiento PCI-DSS          & 0.16 & 16.0\% \\
8  & Integración legacy            & 0.12 & 12.0\% \\
9  & SLA y disponibilidad          & 0.06 & 6.0\%  \\
10 & Gobernanza regulatoria de IA  & 0.06 & 6.0\%  \\
\midrule
\textbf{Total} & & \textbf{1.00} & \textbf{100\%} \\
\bottomrule
\end{tabular}
\end{table}

\subsection*{B.2 Calificaciones por criterio y puntuaciones ponderadas}

La Tabla \ref{tab:4.7} detalla las calificaciones $r_{fi}$ asignadas a cada
framework en escala $[0,10]$ y la contribución ponderada $w_i \cdot r_{fi}$ de
cada criterio a la puntuación global $S_f$.

\begin{table}[htbp]
\centering
\caption{Calificaciones y contribuciones ponderadas por criterio}
\label{tab:4.7}
\scriptsize
\begin{tabular}{@{}clc cc cc cc@{}}
\toprule
& & & \multicolumn{2}{c}{Angular} & \multicolumn{2}{c}{React} & \multicolumn{2}{c}{Vue.js} \\
\cmidrule(lr){4-5}\cmidrule(lr){6-7}\cmidrule(lr){8-9}
N.\textsuperscript{o} & Criterio & $w_i$ & $r_{fi}$ & $w_i r_{fi}$ & $r_{fi}$ & $w_i r_{fi}$ & $r_{fi}$ & $w_i r_{fi}$ \\
\midrule
1  & Seguridad intrínseca         & 0.16 & 9.2 & 1.472 & 7.8 & 1.248 & 7.5 & 1.20 \\
2  & Rendimiento                  & 0.13 & 8.5 & 1.105 & 8.8 & 1.144 & 8.0 & 1.04 \\
3  & Mantenibilidad               & 0.13 & 9.0 & 1.170 & 7.5 & 0.975 & 7.8 & 1.014 \\
4  & Calidad del código IA        & 0.08 & 9.0 & 0.720 & 7.0 & 0.560 & 7.5 & 0.60 \\
5  & Accesibilidad WCAG           & 0.06 & 8.5 & 0.510 & 8.0 & 0.480 & 8.2 & 0.492 \\
6  & Compat. cross-browser        & 0.04 & 8.0 & 0.320 & 8.5 & 0.340 & 8.0 & 0.32 \\
7  & Cumplimiento PCI-DSS         & 0.16 & 9.0 & 1.440 & 7.5 & 1.200 & 7.0 & 1.12 \\
8  & Integración legacy           & 0.12 & 8.8 & 1.056 & 7.8 & 0.936 & 7.2 & 0.864 \\
9  & SLA y disponibilidad         & 0.06 & 8.7 & 0.522 & 8.2 & 0.492 & 7.8 & 0.468 \\
10 & Gobernanza de IA             & 0.06 & 8.8 & 0.528 & 7.6 & 0.456 & 7.3 & 0.438 \\
\midrule
\multicolumn{3}{r}{\textbf{Puntuación global $S_f$}} & \multicolumn{2}{c}{\textbf{8.84}} & \multicolumn{2}{c}{\textbf{7.83}} & \multicolumn{2}{c}{\textbf{7.56}} \\
\bottomrule
\end{tabular}
\end{table}

\noindent \textit{Nota:} los valores $r_{fi}$ fueron asignados por el equipo
técnico del caso de estudio (12 líderes técnicos y 3 oficiales de
cumplimiento) mediante evaluación directa de cada framework en el entorno de
la institución bancaria estudiada, incluyendo la revisión de código generado
por asistentes de IA durante la Fase 2. La puntuación global $S_f =
\sum_{i=1}^{10} w_i \cdot r_{fi}$ sitúa a Angular como la opción óptima para
el proceso de modernización de los canales digitales de la institución.

\section*{Anexo C: Tablas de Referencia para el Proceso AHP}

\subsection*{C.1 Escala fundamental de comparación de Saaty}

La Tabla \ref{tab:4.8} presenta la escala de comparación pareada propuesta
por Saaty (1980), empleada en la Fase 2 de la metodología para asignar los
juicios de importancia relativa entre criterios.

\begin{table}[htbp]
\centering
\caption{Escala fundamental de Saaty para comparaciones pareadas}
\label{tab:4.8}
\small
\begin{tabular}{@{}cp{4cm}p{6.5cm}@{}}
\toprule
Valor $a_{ij}$ & Significado & Ejemplo de aplicación \\
\midrule
1         & Igual importancia          & Dos criterios contribuyen de igual manera al objetivo \\
3         & Importancia moderada       & Un criterio es ligeramente más importante que el otro \\
5         & Importancia fuerte         & Un criterio es claramente más importante que el otro \\
7         & Importancia muy fuerte     & Un criterio es dominante sobre el otro \\
9         & Importancia extrema        & La superioridad de un criterio es absoluta y evidente \\
2, 4, 6, 8 & Valores intermedios        & Compromisos entre juicios adyacentes \\
Recíprocos & $a_{ji} = 1/a_{ij}$        & Criterio $j$ comparado con criterio $i$ \\
\bottomrule
\end{tabular}
\end{table}

\subsection*{C.2 Índice de Consistencia Aleatoria ($RI_n$)}

La Tabla \ref{tab:4.9} presenta los valores tabulados del Índice de
Consistencia Aleatoria $RI_n$ para matrices de orden $n$, utilizados en el
cálculo de la Razón de Consistencia $CR = CI/RI_n$, donde $CI = (\lambda_{
\text{máx}} - n)/(n-1)$. Se considera aceptable un modelo AHP cuando $CR
\leq 0.10$ (Saaty, 1980).

\begin{table}[htbp]
\centering
\caption{Índice de Consistencia Aleatoria tabulado $RI_n$}
\label{tab:4.9}
\small
\begin{tabular}{@{}c*{9}{c}@{}}
\toprule
$n$   & 2    & 3    & 4    & 5    & 6    & 7    & 8    & 9    & 10 \\
\midrule
$RI_n$ & 0.00 & 0.58 & 0.90 & 1.12 & 1.24 & 1.32 & 1.41 & 1.45 & 1.49 \\
\bottomrule
\end{tabular}
\end{table}

Para el caso de estudio presentado en esta tesis ($n = 10$ criterios), se
utilizó $RI_{10} = 1.49$, obteniéndose $CR = 0.06$, valor que confirma la
consistencia de los juicios del panel de expertos.

\subsection*{C.3 Plantilla de evaluación por criterio}

La Tabla \ref{tab:4.10} presenta la evaluación completa del caso de estudio
aplicado en la institución bancaria peruana, con los pesos $w_i$
consensuados por el panel Delphi y las calificaciones $r_{fi} \in [0,10]$
asignadas por el Comité Evaluador para cada combinación framework--criterio.
La puntuación global se calcula como $S_f = \sum_{i=1}^{10} w_i \cdot
r_{fi}$. La tabla es replicable en cualquier institución bancaria
sustituyendo los valores de $w_i$ y $r_{fi}$ según su contexto específico.

\begin{table}[htbp]
\centering
\caption{Plantilla de evaluación --- caso de estudio (Angular, React y Vue.js)}
\label{tab:4.10}
\small
\begin{tabular}{@{}clcccc@{}}
\toprule
N.\textsuperscript{o} & Criterio & $w_i$ & Angular & React & Vue.js \\
\midrule
1  & Seguridad intrínseca         & 0.16 & 9.2 & 7.8 & 7.5 \\
2  & Rendimiento de carga         & 0.13 & 8.5 & 8.8 & 8.0 \\
3  & Mantenibilidad               & 0.13 & 9.0 & 7.5 & 7.8 \\
4  & Calidad del código IA        & 0.08 & 9.0 & 7.0 & 7.5 \\
5  & Accesibilidad WCAG           & 0.06 & 8.5 & 8.0 & 8.2 \\
6  & Compat. cross-browser        & 0.04 & 8.0 & 8.5 & 8.0 \\
7  & Cumplimiento PCI-DSS         & 0.16 & 9.0 & 7.5 & 7.0 \\
8  & Integración legacy           & 0.12 & 8.8 & 7.8 & 7.2 \\
9  & SLA y disponibilidad         & 0.06 & 8.7 & 8.2 & 7.8 \\
10 & Gobernanza de IA             & 0.06 & 8.8 & 7.6 & 7.3 \\
\midrule
\textbf{Puntuación global $S_f$} & \textbf{1.00} & \textbf{8.84} & \textbf{7.83} & \textbf{7.56} \\
\bottomrule
\end{tabular}
\end{table}
